\chapter{Perkalian}\label{ch:g3:mult}

Tiga kantong berisi tujuh kelereng: kamu bisa menjumlahkan
$7 + 7 + 7$, tetapi ada cara yang lebih cepat --- perkalian, operasi
untuk penjumlahan berulang. Bab ini menyusun tabel perkalian dari
gambar, lalu mengajarkan cara bersusun.

\section{Apa arti mengalikan}

\begin{definition}[Perkalian]\label{def:g3:mult:def}
\omterm{def:g2:mult:def}{Hasil kali} $3 \times 7$ berarti ``tiga kali tujuh'':
\[
3 \times 7 = 7 + 7 + 7 = 21 .
\]
Dalam gambar, $3 \times 7$ adalah persegi panjang berisi titik: $3$
baris berisi $7$ titik.
\end{definition}

\begin{omfigure}
\begin{tikzpicture}[scale=0.42]
  \foreach \r in {0,1,2}
    \foreach \c in {0,...,6}
      \fill[omDef] (\c,-\r) circle (0.27);
  \node[right, font=\small] at (7.0,-1.0) {$3$ baris berisi $7$:\ \
  $3 \times 7 = 21$};
  \begin{scope}[yshift=-4.7cm]
  \foreach \r in {0,...,6}
    \foreach \c in {0,1,2}
      \fill[omThm] (\c,-\r) circle (0.27);
  \node[right, font=\small] at (7.0,-3.0) {$7$ baris berisi $3$:\ \
  $7 \times 3 = 21$ juga!};
  \end{scope}
\end{tikzpicture}

{\small $21$ titik yang sama, dibilang per baris atau per kolom:
urutan sebuah \omterm{def:g2:mult:def}{hasil kali} tidak berpengaruh. Mengetahui $3 \times 7$
memberikan $7 \times 3$ cuma-cuma --- tabelnya tinggal separuh!}
\end{omfigure}

\begin{proposition}[Aturan yang membantu]\label{prop:g3:mult:rules}
\begin{enumerate}
  \item Urutannya tidak berpengaruh: $a \times b = b \times a$.
  \item Mengalikan dengan $1$ tidak mengubah apa pun; mengalikan
        dengan $0$ memberi $0$.
  \item Mengalikan dengan $10$ menambahkan sebuah \omterm{def:g1:counting:zero}{nol}:
        $34 \times 10 = 340$; dengan $100$, dua \omterm{def:g1:counting:zero}{nol}.
  \item Sebuah tabel dapat disusun ulang dari tetangganya:
        $6 \times 8 = 5 \times 8 + 8 = 40 + 8 = 48$.
\end{enumerate}
\end{proposition}
\admitted

\begin{example}[Persegi tabel perkalian]\label{ex:g3:mult:table}
Semua tabel muat dalam satu persegi. Untuk membaca $6 \times 7$:
baris $6$, kolom $7$.
\begin{center}
\renewcommand{\arraystretch}{1.15}
\begin{tabular}{c|ccccccccc}
$\times$ & $1$ & $2$ & $3$ & $4$ & $5$ & $6$ & $7$ & $8$ & $9$ \\
\hline
$1$ & 1 & 2 & 3 & 4 & 5 & 6 & 7 & 8 & 9 \\
$2$ & 2 & 4 & 6 & 8 & 10 & 12 & 14 & 16 & 18 \\
$3$ & 3 & 6 & 9 & 12 & 15 & 18 & 21 & 24 & 27 \\
$4$ & 4 & 8 & 12 & 16 & 20 & 24 & 28 & 32 & 36 \\
$5$ & 5 & 10 & 15 & 20 & 25 & 30 & 35 & 40 & 45 \\
$6$ & 6 & 12 & 18 & 24 & 30 & 36 & 42 & 48 & 54 \\
$7$ & 7 & 14 & 21 & 28 & 35 & 42 & 49 & 56 & 63 \\
$8$ & 8 & 16 & 24 & 32 & 40 & 48 & 56 & 64 & 72 \\
$9$ & 9 & 18 & 27 & 36 & 45 & 54 & 63 & 72 & 81 \\
\end{tabular}
\end{center}
Persegi itu simetris terhadap diagonalnya --- itulah aturan 1 yang
terlihat oleh mata.
\end{example}

\section{Perkalian bersusun}

\begin{method}[Mengalikan dengan bilangan satu angka]\label{met:g3:mult:column}
Untuk menghitung $234 \times 6$:
\begin{enumerate}
  \item kalikan satuannya: $4 \times 6 = 24$: tulis $4$, simpan $2$;
  \item kalikan puluhannya lalu tambahkan simpanan:
        $3 \times 6 + 2 = 20$: tulis $0$, simpan $2$;
  \item kalikan ratusannya lalu tambahkan simpanan:
        $2 \times 6 + 2 = 14$: tulis $14$.
\end{enumerate}
\[
\begin{array}{r}
2\,3\,4 \\
\times \quad 6 \\
\hline
1\,4\,0\,4
\end{array}
\]
Taksirlah untuk memeriksa: $234$ dekat dengan $200$, dan
$200 \times 6 = 1\,200$: jawaban $1\,404$ masuk akal.
\end{method}

\begin{example}[Memecah untuk mengalikan di kepala]\label{ex:g3:mult:split}
Cara bersusun diam-diam memecah bilangan menurut tempatnya. Di
kepala, lakukanlah hal yang sama:
\[
23 \times 4 = (20 + 3) \times 4 = 20 \times 4 + 3 \times 4
= 80 + 12 = 92 .
\]
Dan dengan tetangga yang ramah: $99 \times 5 = 100 \times 5 - 5 =
495$.
\end{example}

\begin{omfigure}
\begin{tikzpicture}[scale=0.55]
  \fill[omDef!18] (0,0) rectangle (10,4);
  \fill[omThm!18] (10,0) rectangle (11.5,4);
  \draw[thick] (0,0) rectangle (11.5,4);
  \draw[thick] (10,0) -- (10,4);
  \node[font=\small] at (5,2) {$20 \times 4 = 80$};
  \node[font=\small, rotate=90] at (10.75,2) {$3 {\times} 4$};
  \node[below, font=\small] at (5,0) {$20$};
  \node[below, font=\small] at (10.75,0) {$3$};
  \node[left, font=\small] at (0,2) {$4$};
\end{tikzpicture}

{\small Mengapa memecah itu berhasil: persegi panjang $23 \times 4$
yang tersusun dari persegi satuan adalah bagian $20 \times 4$
ditambah bagian $3 \times 4$ --- $80 + 12 = 92$ persegi.}
\end{omfigure}

\section{Perkalian dalam soal cerita}

\begin{example}\label{ex:g3:mult:problem}
Sebuah minibus mengangkut $9$ penumpang. Berapa penumpang yang dapat
diangkut oleh $27$ minibus?
\begin{enumerate}
  \item Ini perkalian: $27$ kelompok berisi $9$.
  \item $27 \times 9 = 243$ (secara bersusun, atau
        $27 \times 9 = 27 \times 10 - 27 = 270 - 27$).
  \item Kalimatnya: minibus itu dapat mengangkut $243$ penumpang.
\end{enumerate}
\end{example}

\section{Latihan}

\begin{exercise}[$\star$]\label{exo:g3:mult:1}
Tulislah sebagai perkalian, lalu hitunglah:
$8 + 8 + 8 + 8$; \quad $5 + 5 + 5$; \quad
$20 + 20 + 20 + 20 + 20$.
\end{exercise}

\begin{exercise}[$\star$]\label{exo:g3:mult:2}
Sebutkan lalu lengkapi: $6 \times 7$; \ $8 \times 4$; \ $9 \times 9$;
\ $7 \times 8$; \ $6 \times 6$.
\end{exercise}

\begin{exercise}[$\star$]\label{exo:g3:mult:3}
Gambarlah persegi panjang titik untuk $4 \times 6$, lalu untuk
$6 \times 4$. Apa yang kamu perhatikan pada kedua hitungan itu?
\end{exercise}

\begin{exercise}[$\star$]\label{exo:g3:mult:4}
Hitunglah dengan aturan pada \cref{prop:g3:mult:rules}:
$56 \times 10$; \quad $8 \times 100$; \quad $73 \times 0$; \quad
$45 \times 1$; \quad $30 \times 10$.
\end{exercise}

\begin{exercise}[$\star$]\label{exo:g3:mult:5}
Hitunglah secara bersusun: $132 \times 3$; \quad $217 \times 4$; \quad
$408 \times 7$.
\end{exercise}

\begin{exercise}[$\star$]\label{exo:g3:mult:6}
Hitunglah di kepala dengan memecah (seperti pada \cref{ex:g3:mult:split}):
$31 \times 5$; \quad $42 \times 3$; \quad $25 \times 6$.
\end{exercise}

\begin{exercise}[$\star$]\label{exo:g3:mult:7}
Hitunglah dengan tetangga yang ramah: $99 \times 7$; \quad
$101 \times 6$; \quad $98 \times 5$.
\end{exercise}

\begin{exercise}[$\star$]\label{exo:g3:mult:8}
Taksirlah dulu, lalu hitunglah: $389 \times 5$ (taksir dengan $400$);
$612 \times 8$ (taksir dengan $600$).
\end{exercise}

\begin{exercise}[$\star$]\label{exo:g3:mult:9}
Satu kotak telur berisi $6$ telur. Berapa telur di dalam $38$ kotak?
Tulislah operasinya, hitunglah, lalu jawablah dengan satu kalimat.
\end{exercise}

\begin{exercise}[$\star\star$]\label{exo:g3:mult:10}
Sebuah ruang kelas punya $8$ baris berisi $4$ meja; setiap meja
ditempati $2$ murid. Berapa murid yang dapat ditampung ruangan itu? (Dua perkalian.)
\end{exercise}

\begin{exercise}[$\star\star$]\label{exo:g3:mult:11}
Carilah semua cara menyusun $24$ kursi menjadi baris yang sama
banyak (baris berisi $1$, berisi $2$, \dots). Susunan mana saja yang
mungkin? (Pakailah persegi tabel: carilah $24$ di dalamnya.)
\end{exercise}

\begin{exercise}[$\star\star$]\label{exo:g3:mult:12}
Zahra menghitung $47 \times 6$ begini: $40 \times 6 = 240$,
$7 \times 6 = 42$, lalu $240 + 42 = 282$. Lia menghitung
$47 \times 6 = 50 \times 6 - 3 \times 6 = 300 - 18 = 282$. Jelaskan
kedua cara itu. Mana yang kamu sukai untuk $38 \times 4$? Hitunglah
dengan kedua cara.

\end{exercise}
